\section{Síntesis y ampliación}

\subsection{Repaso de los puntos clave}

Esta clase gira en torno a "cómo resolver eficientemente la ODE inversa de los modelos de difusión", con las siguientes contribuciones principales:

\begin{enumerate}
    \item \textbf{Descubrimiento de una estructura semilineal}: El término derecho de la PF-ODE se puede descomponer en una parte lineal (que se puede resolver exactamente) y una parte no lineal (que requiere aproximación numérica), lo que sienta las bases para una resolución eficiente.
    \item \textbf{Aplicación de integradores exponenciales}: Mediante el método del factor integrante, se trata la parte lineal con precisión, limitando la aproximación numérica únicamente a la integración de la predicción del ruido.
    \item \textbf{Parametrización en el espacio $\lambda$}: Realizando un cambio de variable con log-SNR, se simplifica enormemente la expresión de la integral.
    \item \textbf{Reinterpretación de DDIM}: DDIM es un caso particular de primer orden de DPM-Solver, no del método de Euler ingenuo.
    \item \textbf{Construcción recursiva de métodos de alto orden}: Realizando una aproximación polinómica de alto orden a la predicción del ruido en el espacio $\lambda$, se obtiene un aumento significativo de precisión a cambio de pocas NFE adicionales.
    \item \textbf{DPM-Solver++}: Se convierte la aproximación polinómica para realizar predicciones de datos, mejorando aún más la precisión con pasos limitados.
\end{enumerate}

\subsection{Una perspectiva más amplia}

\begin{knowledgebox}{Por qué es tan importante la comprensión matemática}
El ponente enfatizó repetidamente en el curso que todas estas aceleraciones se logran bajo el prerrequisito de \textbf{no modificar el modelo entrenado}. Esto refleja profundamente un principio: la comprensión matemática de la estructura del problema puede traducirse directamente en mejoras de rendimiento práctico. Desde DDPM hasta DDIM y luego DPM-Solver, cada paso representa una comprensión más profunda del mismo proceso de difusión.
\end{knowledgebox}

\subsection{Desarrollos futuros}

Tras DPM-Solver, la aceleración del muestreo en modelos de difusión ha seguido evolucionando rápidamente:

\begin{itemize}
    \item \textbf{Modelos de consistencia} (Consistency Models, Song et al., 2023): Aprenden un mapeo directo desde el ruido hasta los datos, permitiendo la generación en un solo paso.
    \item \textbf{Destilación progresiva} (Progressive Distillation, Salimans \& Ho, 2022): Comprime procesos de múltiples pasos en pocos pasos mediante destilación de conocimiento.
    \item \textbf{Flow Matching} (Lipman et al., 2022): Diseña trayectorias de transporte más rectas, permitiendo que el método de Euler simple funcione de manera eficiente.
    \item \textbf{Rectified Flow} (Liu et al., 2022): Hace que las trayectorias de la ODE tiendan a ser rectas mediante operaciones de reflujo.
\end{itemize}

\subsection{Lecturas adicionales}

\begin{itemize}
    \item Lu, C., et al. (2022). \textit{DPM-Solver: A Fast ODE Solver for Diffusion Probabilistic Model Sampling in Around 10 Steps}. NeurIPS 2022.
    \item Lu, C., et al. (2022). \textit{DPM-Solver++: Fast Solver for Guided Sampling of Diffusion Probabilistic Models}. arXiv:2211.01095.
    \item Song, J., et al. (2020). \textit{Denoising Diffusion Implicit Models}. ICLR 2021.
    \item Song, Y., et al. (2021). \textit{Score-Based Generative Modeling through Stochastic Differential Equations}. ICLR 2021.
    \item Karras, T., et al. (2022). \textit{Elucidating the Design Space of Diffusion-Based Generative Models}. NeurIPS 2022.
\end{itemize}

%% --- Fin del contenido --- %%

\end{document}
